\section{\dandelion's Approach}

To achieve high elasticity while eliminating the need to pre-provision idle sandboxes in memory, we propose \dandelion.
Our key insight is that redesigning the traditional POSIX-based programming interface for cloud applications in the cloud-native era enables co-designing a more efficient and elastic cloud platform. We observe that cloud-native applications typically consist of interactions between various cloud services (exposed over REST APIs) and snippets of user-defined compute logic. Hence, we propose a programming model in which applications are expressed as DAGs of \textit{pure compute functions} and \textit{communication functions}.
Pure compute functions are snippets of untrusted user code that consume a declared set of inputs and produce a declared set of outputs.
Communication functions enable interaction with external services.
They are implemented by the \dandelion platform. Users invoke communication functions from user application DAGs, but for security, the implementation of communication functions cannot be modified by users. Currently, \dandelion supports an HTTP communication function that enables interaction with remote services like cloud storage buckets, AI inference services, and data warehouses over REST APIs.  
Writing a \dandelion application involves: 1) expressing the DAG of compute and communication functions using a domain-specific language (DSL) and 2) providing the binaries for user-defined compute functions. \S\ref{sec:design-app} describes the DSL for expressing application DAGs (which also makes function inputs and outputs explicit),  language support for developing user-defined compute functions, and the library of currently supported communication functions.
This programming model is inspired by prior dataflow systems~\cite{spark, beam, naiad, dryad, dryadLINQ, malte-dataflow}. Similar to systems like Spark~\cite{spark}, \dandelion's programming model makes data dependencies and thus data-parallelism and task-parallelism explicit to the platform. In contrast with traditional dataflow graphs, where edges encode internal communication and external communication is only performed at the dataflow sources and sinks, \dandelion function instantiations can either perform local compute or exchange data with other cloud services. Furthermore, we design \dandelion to execute dataflow DAGs with secure isolation in multi-tenant deployments without relying on execution in traditional VMs, which are slow to boot.

Decomposing a cloud application into pure compute and communication functions allows \dandelion to optimize the execution system for high elasticity while maintaining secure isolation guarantees (\S\ref{sec:design:system-arch}).
It enables faster cold starts: pure compute functions can be securely executed in lightweight sandboxes that provide memory isolation (not necessarily VMs), since they do not rely on extra software environments like a guest OS.
The separation of compute and communication also enables \dandelion to use cooperative network I/O and elastically adjust the CPU core allocation based on the number of compute vs. communication functions in the system. The platform can monitor each application's I/O intensity based on its number of compute vs. communication functions. \dandelion leverages this information for CPU core allocation, over-committing CPU cores for communication functions and dedicating CPU cores to compute functions.

\textbf{Target applications.} We design \dandelion for workloads with spiky load patterns (scaling from 0/few to 1/many with long periods of no/low load) or data-dependent parallelism. Many applications can benefit from \dandelion's ability to quickly scale-out to multiple CPU cores when an application computes on many independent inputs and quickly scale-in when the application needs to process data sequentially.
Our goal is to enable fine-grain~\cite{granular-computing}, elastic computing with secure isolation for a variety of cloud-native applications, such as data transformation pipelines (e.g., image compression, video transcoding~\cite{opencv, excamera}), distributed log processing~\cite{minio-log-processing}, CI/CD tasks, elastic query processing (which can include untrusted user-defined function (UDF) execution~\cite{containerized-udf}), and emerging agentic AI workflows~\cite{compound-ai-blog, agentic-ai-databricks}.

\textbf{Non-goals.} However, \dandelion is not designed for applications that require large, frequent state synchronization, like online transaction processing, online gaming services, AI training, or data processing algorithms that rely on fine-grain multithreading and shared memory for performance.